% ============================================================================
%  CO CHE SUP DO CUA CAC TAN CONG KHOI PHUC NHAN DUNG THONG KE PHU TRO
%
%  Ban RUT GON: chi giu chuoi suy luan CAN VA DU de chung minh su sup do cua
%  nhom dung thong ke phu tro (ZLG+, RLU).
%
%  Ban DAY DU (co doi chieu LLG/LLG+/VLI*, bat bien thang do, phep thu thuoc
%  tinh ham, phan "thong tin nam o huong") nam o:
%      VLI_cochesupdo_vi_full.tex
%
%  File NAY DOC LAP voi VLI_baocao_vi.tex (khong ghi de, khong thay the).
%
%  NGUYEN TAC: khong trich mot con so mo phong nao. Chi gom dai so chung minh
%  duoc va cac su that doc truc tiep tu ma nguon (tep + ham ghi o Phu luc).
%
%  Compile:
%     pdflatex VLI_cochesupdo_vi.tex   (chay 3 lan)
% ============================================================================
\documentclass[11pt,a4paper]{article}

\usepackage{iftex}
\ifPDFTeX
  \usepackage[utf8]{inputenc}
  \usepackage[T5]{fontenc}
\else
  \usepackage{fontspec}
\fi
\usepackage[vietnamese]{babel}

\usepackage[margin=2.3cm]{geometry}
\usepackage{amsmath,amssymb,amsthm,mathtools}
\usepackage{booktabs}
\usepackage{array}
\usepackage[hidelinks]{hyperref}

\newtheorem{dinhly}{Định lý}
\newtheorem{bode}[dinhly]{Bổ đề}
\newtheorem{menhde}[dinhly]{Mệnh đề}
\newtheorem{hequa}[dinhly]{Hệ quả}
\newtheorem{dinhnghia}[dinhly]{Định nghĩa}
\theoremstyle{remark}
\newtheorem{nhanxet}[dinhly]{Nhận xét}
\renewcommand{\proofname}{Chứng minh}

\newcommand{\R}{\mathbb{R}}
\newcommand{\one}{\mathbf{1}}
\newcommand{\pos}[1]{\left[#1\right]_+}
\newcommand{\quant}{\mathcal{Q}}
\newcommand{\epstr}{\bar\varepsilon_{\mathrm{tr}}}
\newcommand{\epsaux}{\bar\varepsilon_{\mathrm{aux}}}
\newcommand{\code}[1]{\texttt{\small #1}}

\renewcommand{\arraystretch}{1.25}

\title{Cơ chế sụp đổ của các tấn công khôi phục nhãn\\[0.3em]
\large dùng thống kê phụ trợ}
\author{}
\date{\today}

\begin{document}
\maketitle

\begin{abstract}
Báo cáo này chứng minh một mệnh đề duy nhất: khi mô hình toàn cục được huấn
luyện tốt, đầu ra của ZLG$^+$ và RLU bị \emph{nhốt} trong một quả cầu bán
kính $A+C$ quanh tiên nghiệm lớp của tập phụ trợ --- một vectơ không chứa
thông tin nào về nhãn của nạn nhân --- với $A=\Theta(\epstr)$. Lập luận gồm
ba mắt xích: (i) độ lớn gradient lớp cuối \emph{chính là} sai số huấn luyện,
nên tử số co; (ii) thước mà hai phương pháp này so vào lấy từ tập phụ trợ nên
mẫu số không co; (iii) một chặn đi dọc đoạn thẳng cho thấy tỉ số $A$ của hai
đại lượng đó chặn trực tiếp khoảng cách từ đầu ra tới tâm. Mọi khẳng định về
mã nguồn đều ghi rõ tệp và hàm; mọi đại lượng chưa đo đều ghi rõ là
\textbf{chưa đo}.
\end{abstract}

\tableofcontents

% ===========================================================================
\section{Phạm vi}
\label{sec:pham-vi}

\subsection{Mục tiêu và những gì đã cắt}

Bản này chỉ giữ chuỗi suy luận \emph{cần và đủ} để chứng minh sự sụp đổ của
nhóm dùng thống kê phụ trợ. Đã cắt khỏi bản đầy đủ
(\code{VLI\_cochesupdo\_vi\_full.tex}, vẫn nằm cùng thư mục):

\begin{itemize}
\item phần đối chiếu với LLG, LLG$^+$, VLI$^*$ và khái niệm \emph{bất biến
  thang độ}: đó là câu trả lời cho ``vì sao nhóm kia \emph{không} sụp'', một
  câu hỏi khác;
\item phép thử thuộc tính hàm $f(\lambda\Gamma)$ so với $f(\Gamma)$: cùng lý
  do trên;
\item phần ``thông tin nằm ở hướng'' (giả nghịch đảo, tái lập nhãn): nó chứng
  minh sụp đổ \emph{tránh được}, không chứng minh sụp đổ \emph{xảy ra};
\item giả thiết không-kẹp và đồng nhất thức affine đi kèm: từ khi Mệnh đề
  \ref{md:bay} được chứng minh vô điều kiện, chúng không còn nằm trên đường
  tới kết luận.
\end{itemize}

\subsection{Quy tắc trích dẫn}

Mọi khẳng định mang nhãn \textsc{[cm]}: chứng minh đầy đủ từ định nghĩa, hoặc
đọc trực tiếp từ mã nguồn (tệp và hàm ghi ở Phụ lục \ref{sec:phuluc}). Không
có con số mô phỏng nào trong bản này. Đại lượng chưa đo trên pipeline thật
được ghi \textbf{chưa đo}, và không có số nào thay cho nó.

\subsection{Mệnh đề cần chứng minh}

\begin{quote}
Khi mô hình được huấn luyện tốt, đầu ra của ZLG$^+$ và RLU hội tụ về một
vectơ không phụ thuộc nhãn của client --- cụ thể là tiên nghiệm lớp của tập
phụ trợ của chính kẻ tấn công.
\end{quote}

% ===========================================================================
\section{Ký hiệu}
\label{sec:ki-hieu}

Máy chủ trung thực-nhưng-tò mò quan sát cập nhật của một client sau cửa sổ
gồm $m$ bước SGD cục bộ, mỗi bước trên một lô $B$ mẫu; tổng số nhãn $n=mB$.
Cần khôi phục $N=(N_1,\dots,N_C)$ với $\sum_j N_j=n$.

\begin{center}
\begin{tabular}{@{}cl@{}}
\toprule
\textbf{Ký hiệu} & \textbf{Nghĩa} \\
\midrule
$C$, $B$, $m$, $n=mB$ & số lớp; cỡ lô; số bước cục bộ; tổng số nhãn \\
$y_i$, $e_{y_i}$ & nhãn thật của mẫu $i$; vectơ one-hot của nó \\
$p_i$ & phân phối softmax mô hình gán cho mẫu $i$ \\
$\varepsilon_i=1-p_{i,y_i}$ & \emph{sai số} của mẫu $i$; $\varepsilon_i\in[0,1)$ \\
$N_j$, $\bar\varepsilon_j$ & số mẫu lớp $j$; sai số trung bình của lớp $j$ \\
$a_j=N_j\bar\varepsilon_j$ & \emph{khối lượng lỗi} của lớp $j$ \\
$\epstr=n^{-1}\sum_i\varepsilon_i$ & sai số huấn luyện trung bình trên cửa sổ \\
$\epsaux$ & sai số trung bình trên tập phụ trợ của kẻ tấn công \\
$\bar p_j=1-\bar\varepsilon_{\mathrm{aux},j}$ & độ tự tin trên \emph{lớp đúng} $j$, đo trên phụ trợ \\
$s_i=\sum_d e_{i,d}$, $\bar O$ & tổng đặc trưng của mẫu $i$ ($\ge0$); trung bình của $s_i$ trên phụ trợ \\
$\Gamma$ & tín hiệu gradient mà phương pháp đọc được \\
$\quant$ & phép làm tròn về số nguyên \\
\bottomrule
\end{tabular}
\end{center}

Hai dạng tín hiệu xuất hiện trong mã nguồn:
\begin{align}
G &= \sum_{i=1}^{n}\bigl(p_i-e_{y_i}\bigr) && \text{(gradient bias; RLU)},
\label{eq:G}\\
\tilde G_j &= \sum_{i=1}^{n} w_i\bigl(p_{i,j}-\delta_{j,y_i}\bigr),
 \quad w_i=\frac{s_i}{\bar O} && \text{(gradient trọng số; ZLG$^+$)}.
\label{eq:Gtilde}
\end{align}

Ma trận nhầm lẫn chuẩn hóa, với $Q_{jj}:=0$:
\begin{equation}
Q_{jk}\;=\;\frac{1}{a_j}\sum_{i:\,y_i=j}p_{i,k}\qquad(k\neq j).
\label{eq:Q}
\end{equation}
Phép chia cho $a_j$ khử $\bar\varepsilon$ khỏi $Q$: tử và mẫu cùng thứ nguyên
$[\varepsilon]^1$. Nên $Q$ mô tả \emph{hình dạng} của lỗi, không mô tả
\emph{độ lớn}. Phân biệt này là bản lề của cả báo cáo.

% ===========================================================================
\section{Mắt xích 1: tử số co theo sai số}
\label{sec:tu-so}

\begin{bode}[$Q$ ngẫu nhiên theo hàng]
\label{bd:Qstoch}
\textsc{[cm]}\; Với mọi $j$ sao cho $a_j>0$: $\sum_k Q_{jk}=1$.
\end{bode}

\begin{proof}
$\sum_{k\ne j}\sum_{i:y_i=j}p_{i,k}=\sum_{i:y_i=j}(1-p_{i,j})
=\sum_{i:y_i=j}\varepsilon_i=a_j$; chia hai vế cho $a_j$.
\end{proof}

\begin{bode}[Đồng nhất thức tín hiệu]
\label{bd:tinhieu}
\textsc{[cm]}\; Với $a=(a_1,\dots,a_C)$: $\;G=-\bigl(I-Q^{\!\top}\bigr)a$.
\end{bode}

\begin{proof}
$G_j=\sum_i p_{i,j}-N_j$. Tách theo nhãn thật:
$\sum_{i:y_i=j}p_{i,j}=N_j-a_j$, và $\sum_{i:y_i=k}p_{i,j}=a_kQ_{kj}$ với
$k\ne j$. Dùng $Q_{jj}=0$ để gộp $k=j$ vào tổng mà không đổi giá trị:
$G_j=(N_j-a_j)+\sum_k a_kQ_{kj}-N_j=-\bigl[(I-Q^{\!\top})a\bigr]_j$.
\end{proof}

\begin{bode}[Độ lớn tín hiệu bị chặn bởi sai số huấn luyện]
\label{bd:dolon}
\textsc{[cm]}\;
\begin{equation}
\|G\|_1\;\le\;2n\,\epstr,
\qquad
\|\tilde G\|_1\;\le\;2n\,w_{\max}\,\epstr .
\label{eq:dolon}
\end{equation}
\end{bode}

\begin{proof}
Từ Bổ đề \ref{bd:tinhieu} và Bổ đề \ref{bd:Qstoch}:
$\|G\|_1\le\|a\|_1+\|Q^{\!\top}a\|_1=2\|a\|_1$, và
$\|a\|_1=\sum_j N_j\bar\varepsilon_j=\sum_i\varepsilon_i=n\epstr$. Với
$\tilde G$: mỗi mẫu góp $w_i\|p_i-e_{y_i}\|_1=2w_i\varepsilon_i\le
2w_{\max}\varepsilon_i$.
\end{proof}

\begin{bode}[Tổng triệt tiêu]
\label{bd:tongkhong}
\textsc{[cm]}\; $\sum_{j=1}^{C}\Gamma_j=0$ với $\Gamma=G$ \emph{và} với
$\Gamma=\tilde G$, bất kể trọng số $w_i$.
\end{bode}

\begin{proof}
Với mỗi mẫu $i$: $\sum_j(p_{i,j}-\delta_{j,y_i})=1-1=0$, vì $p_i$ là phân
phối xác suất và $e_{y_i}$ là one-hot. Mọi tổ hợp tuyến tính theo $i$ giữ
tính chất này.
\end{proof}

\begin{nhanxet}[Đây là chỗ không có đường thoát]
\label{nx:khongthoat}
Bổ đề \ref{bd:dolon} không phải hiện tượng phụ của kiến trúc hay dữ liệu.
Gradient của lớp tuyến tính cuối dưới entropy chéo là \emph{tổng các phần dư}
$p_i-e_{y_i}$, và chuẩn $L_1$ của một phần dư bằng đúng $2\varepsilon_i$. Nên
độ lớn của tín hiệu \emph{chính là} sai số, không phải chỉ tương quan với nó.
Huấn luyện mô hình tốt lên \emph{bắt buộc} làm $\|\Gamma\|_1$ co. Mọi tấn
công đều phải sống với điều này; câu hỏi duy nhất là nó có \emph{quan tâm}
hay không.

Giả thiết duy nhất ở đây là $w_{\max}=\max_i s_i/\bar O$ bị chặn, và nó chỉ
cần cho $\tilde G$ (tức ZLG$^+$). Giá trị $w_{\max}$ thực tế: \textbf{chưa
đo}.
\end{nhanxet}

% ===========================================================================
\section{Đối tượng: dạng chuẩn đọc từ mã nguồn}
\label{sec:dang-chuan}

\begin{menhde}[Dạng chuẩn]
\label{md:dangchuan}
\textsc{[cm]}\; ZLG$^+$ và bước đầu của RLU đều tính
\begin{equation}
\hat N\;=\;\quant\Bigl(R\bigl(c-\beta\Gamma\bigr)\Bigr),
\qquad R(x)=\frac{n\,\pos{x}}{\|\pos{x}\|_1},
\label{eq:dangchuan}
\end{equation}
với $c\in\R^C_{\ge0}$ \emph{không phụ thuộc nhãn của cửa sổ bị tấn công}. Cụ
thể: ZLG$^+$ có $c=n\bar p$, $\beta=1$, $\Gamma=\tilde G$; RLU có
$c=\bar\rho\cdot\mathrm{softmax}(\bar\mu)$, $\beta=1/n$, $\Gamma=G$.
\end{menhde}

\begin{proof}[Chứng minh (đọc mã nguồn)]
\emph{ZLG$^+$} (\code{client/attacks.py}, \code{ZLGp\_attack}). Mã tính
\[
n_j\;=\;mB\Bigl(\bar p_j-\frac{\mathrm{grad}_j}{\bar O}\Bigr),
\qquad
\mathrm{grad}_j=\sum_d\Delta W_{jd}=\frac{\bar O\,\tilde G_j}{n},
\]
nên $n_j=n\bar p_j-\tilde G_j$; rồi cắt âm, chuẩn hóa về tổng $n$, làm tròn.

\emph{RLU} (\code{client/client\_utils.py},
\code{estimated\_entropy\_from\_grad}). Mã đọc nguyên văn là
\begin{center}
\begin{tabular}{@{}ll@{}}
\toprule
\code{bias = -np.array(bias)}             & $-G/n$ \\
\code{bias[i] = bias[i] + shift[i]}       & $x=c-G/n$ \\
\code{if bias[i] < 0: bias[i] = 0}        & $\pos{x}$ \\
\code{s = np.sum(abs(bias))}              & $\|\pos{x}\|_1$ \\
\code{solution[i] = round(bias[i]/s * B)} & $\quant\bigl(n\pos{x}/\|\pos{x}\|_1\bigr)$ \\
\bottomrule
\end{tabular}
\end{center}
với \code{shift = softmax(new\_mu) * rho\_mean}.
\end{proof}

\begin{nhanxet}[Phạm vi: RLU chạy hai vòng]
\label{nx:haivong}
\textsc{[cm]}\; Hàm \code{RLU\_attack} chạy dạng chuẩn \emph{hai lần}; giá
trị nó trả về là của \textbf{vòng hai}. Thước của vòng hai do
\begin{center}
\code{estimate\_static\_RLU\_with\_posterior(args, n, ...)}
\end{center}
sinh ra --- mà tham số \code{n} của nó chính là ước lượng của vòng một, và
được dùng trực tiếp trong
\begin{center}
\code{gb[i] = -n[i]/args.batch\_size + softmax(...)[i]}
\end{center}

Hệ quả: \textbf{tâm của vòng hai phụ thuộc $\Gamma$}, tức phụ thuộc nhãn.
Mọi kết quả dưới đây được chứng minh cho \emph{vòng một} của RLU. Phần còn
giữ nguyên cho vòng hai là $\|c\|_1=\bar\rho$ (vì softmax vẫn tổng bằng $1$),
nên Mệnh đề \ref{md:bay} vẫn cho tính bị nhốt --- nhưng quanh một tâm không
còn vô can. Xem chỗ không phủ số \ref{item:hvong}.
\end{nhanxet}

\begin{nhanxet}[Tâm của ánh xạ]
\label{nx:tam}
Nếu $\Gamma=0$ thì \eqref{eq:dangchuan} cho $\quant\bigl(R(c)\bigr)$, một
vectơ phụ thuộc \emph{duy nhất} vào $c$, tức vào tập phụ trợ. Gọi nó là
\textbf{tâm}. Với ZLG$^+$ tâm là $n\bar p/\sum_j\bar p_j$, với RLU là
$n\cdot\mathrm{softmax}(\bar\mu)$ --- cả hai đều là tiên nghiệm lớp của dữ
liệu phụ trợ của kẻ tấn công. Toàn bộ phần còn lại của báo cáo chứng minh
rằng $\hat N$ không đi xa được khỏi tâm này.
\end{nhanxet}

\begin{nhanxet}[$\bar O$ không triệt tiêu]
\label{nx:Obar}
Trong chứng minh trên, $\bar O$ biến mất ở bước cuối. Đó là \emph{kế toán ký
hiệu}: định nghĩa \eqref{eq:Gtilde} của $\tilde G$ đã chứa $\bar O$ ở mẫu số
của $w_i$. Viết theo đúng đại lượng máy chủ \emph{quan sát được}:
$\beta=n/\bar O$ và $\Gamma=\mathrm{grad}$. Nên $\bar O$ là một thành phần
nữa của thước --- ước lượng $\bar O$ càng lớn thì tín hiệu càng bị nén. Hành
vi của $\bar O$ khi huấn luyện: \textbf{chưa đo}, và không suy ra được bằng
đại số.
\end{nhanxet}

% ===========================================================================
\section{Mắt xích 2: mẫu số không co}
\label{sec:mau-so}

\begin{menhde}[Chuẩn của thước bị ghim]
\label{md:ghim}
\textsc{[cm]}\;
\begin{enumerate}
\item \textbf{ZLG$^+$}: $\|c\|_1=n\sum_j\bar p_j=nC\bigl(1-\epsaux\bigr)$.
  Đại lượng này \emph{tăng} về trần $nC$ khi tập phụ trợ được phân loại tốt
  hơn.
\item \textbf{RLU}: $\|c\|_1=\bar\rho$, vì $\mathrm{softmax}(\bar\mu)$ tổng
  bằng $1$ với mọi $\bar\mu$. Hơn nữa \code{rlu\_rho\_mean(args)} chỉ đọc
  \code{local\_epochs}, \code{gamma}, \code{lr}, \code{mu},
  \code{fl\_scheme}; nó \emph{không} chạm vào mô hình lẫn tập phụ trợ, nên
  $\bar\rho$ là \textbf{hằng số qua huấn luyện}. Với $\gamma=0{,}1$, $m=10$
  thì $\bar\rho\approx1{,}0988$.
\end{enumerate}
\end{menhde}

\begin{proof}
(1) $\bar p_j=1-\bar\varepsilon_{\mathrm{aux},j}$ (đọc
\code{estimate\_static\_ZLG}: tích lũy \code{pj[k] += mean\_probs[k]}, tức độ
tự tin trên lớp \emph{đúng}), nên
$\sum_j\bar p_j=C-\sum_j\bar\varepsilon_{\mathrm{aux},j}=C(1-\epsaux)$.
(2) Đọc \code{rlu\_rho\_mean} trong \code{client/attacks.py}; công thức
$\rho_t=(1-\gamma^{\,m-t})/(1-\gamma)$ chỉ chứa hyperparameter.
\end{proof}

\begin{nhanxet}[Trả lời phản biện: phụ trợ cũng được huấn luyện tốt lên]
\label{nx:phanbien}
Phản biện tự nhiên: tập phụ trợ cũng chạy qua cùng mô hình, nên $\epsaux$
cũng giảm --- sao thước lại không co?

Vì tiêu chí không phải ``có dùng phụ trợ hay không'' mà là \emph{đại lượng đó
neo ở đâu}. Thước ở đây không phải $\epsaux$ mà là $1-\epsaux$, và nó neo ở
$1$: khi mô hình tốt lên, nó \textbf{tăng} chứ không giảm. Nên $\epsaux$ giảm
càng làm $A$ nhỏ thêm --- phản biện đi ngược lại chính nó. Với RLU còn mạnh
hơn: thước hoàn toàn không đụng tới dữ liệu (Mệnh đề \ref{md:ghim}.2), nên
không có đường nào để nó co.
\end{nhanxet}

\begin{dinhnghia}[Tỉ số tín hiệu trên thước]
\label{dn:A}
\begin{equation}
A\;:=\;\frac{n\,\beta\,\|\Gamma\|_1}{\|c\|_1}\;.
\label{eq:A}
\end{equation}
\end{dinhnghia}

\begin{hequa}
\label{hq:Aeps}
\textsc{[cm]}\; $A=\Theta(\epstr)$: tử số là $\Theta(n\epstr)$ (Bổ đề
\ref{bd:dolon}), mẫu số bị ghim (Mệnh đề \ref{md:ghim}). Với ZLG$^+$ mẫu số
còn \emph{tăng} khi $\epsaux$ giảm, nên $A$ bị bóp từ cả hai phía.
\end{hequa}

% ===========================================================================
\section{Mắt xích 3: tỉ số nhỏ nhốt đầu ra}
\label{sec:bay}

\begin{bode}[Phân rã phần dương]
\label{bd:phanra}
\textsc{[cm]}\; Với mọi $x\in\R^C$:
$\ \|\pos{x}\|_1=\sum_j x_j+\sum_{j:x_j<0}|x_j|$.
\end{bode}

\begin{proof}
$[t]_+=t+[-t]_+$ với mọi $t\in\R$; lấy tổng theo $j$.
\end{proof}

\begin{menhde}[Chặn bẫy]
\label{md:bay}
\textsc{[cm]}\; Giả sử $c\ge0$, $c\ne0$ (đúng cho cả hai phương pháp theo
Mệnh đề \ref{md:dangchuan}). Khi đó
\begin{equation}
\bigl\|R(c-\beta\Gamma)-R(c)\bigr\|_1\;\le\;A ,
\label{eq:dichchuyen}
\end{equation}
và do đó
\begin{equation}
\bigl\|\,\hat N-\underbrace{\quant\bigl(R(c)\bigr)}_{\text{tâm}}\,\bigr\|_1
\;\le\;A+C .
\label{eq:bay}
\end{equation}
Không cần giả thiết nào khác --- nói riêng, \emph{không} cần giả định các tọa
độ không bị cắt.
\end{menhde}

\begin{proof}
Tham số $\tau$ dưới đây là \emph{công cụ của chứng minh}, không phải đại
lượng của bài toán: nó chạy dọc đoạn thẳng nối tâm với điểm thật, với
$\tau=0$ cho $c$ và $\tau=1$ cho $c-\beta\Gamma$, và nó biến mất sau bước lấy
tích phân. (Nó không liên quan tới chỉ số bước cục bộ $t$ của $\rho_t$ trong
Mệnh đề \ref{md:ghim}.)

Đặt $g=\beta\Gamma$, $L=\|c\|_1$ và $\varphi(\tau)=R(c-\tau g)$ với
$\tau\in[0,1]$.
Ký hiệu $S(\tau)=\{j:c_j-\tau g_j>0\}$, $g_S$ là $g$ thu hẹp về $S$, và
$D(\tau)=\|\pos{c-\tau g}\|_1$. Ý tưởng: \emph{đi dọc đoạn thẳng từ tâm tới điểm
thật và chặn tốc độ}; quãng đường không vượt quá tốc độ tối đa nhân thời
gian, mà thời gian bằng $1$.

\emph{Bước 1 ($D\ge L$).} Theo Bổ đề \ref{bd:tongkhong},
$\sum_j(c_j-\tau g_j)=\sum_j c_j=L$. Theo Bổ đề \ref{bd:phanra},
$D(\tau)=L+\sum_{j\notin S(\tau)}|c_j-\tau g_j|\ge L>0$. Nói riêng $\varphi$ xác định
trên cả đoạn, và việc bị cắt chỉ làm mẫu số \emph{tăng}.

\emph{Bước 2 (dấu ngoài $S$).} Nếu $j\notin S(\tau)$ thì $c_j-\tau g_j\le0$; vì
$c_j\ge0$ nên $\tau g_j\ge c_j\ge0$. Vậy \emph{mọi tọa độ bị cắt đều có
$g_j\ge0$}.

\emph{Bước 3 (tốc độ).} Mỗi $c_j-\tau g_j$ affine theo $\tau$ nên đổi dấu nhiều nhất
một lần; $S(\tau)$ chỉ thay đổi tại hữu hạn điểm, và $\varphi$ liên tục,
Lipschitz từng khúc. Trên một khúc, đặt $\sigma=\sum_{j\in S}g_j$ thì
$D'=-\sigma$ và
\[
\varphi'\;=\;-\frac{n}{D}\,g_S+\frac{\sigma}{D}\,\varphi,
\qquad\text{nên}\qquad
\|\varphi'\|_1\;\le\;\frac{n}{D}\Bigl(\|g_S\|_1+|\sigma|\Bigr),
\]
dùng $\|\varphi\|_1=n$ (ảnh của $R$ luôn nằm trên đơn hình khối lượng $n$).
Hai số hạng ứng với hai lực: vectơ tự dịch chuyển, và mẫu số chuẩn hóa thay
đổi.

\emph{Bước 4 (hai mảnh khớp khít).} Theo Bổ đề \ref{bd:tongkhong},
$\sigma=-\sum_{j\notin S}g_j$; theo Bước 2 mọi số hạng đó không âm, nên
$|\sigma|=\|g_{S^c}\|_1$ --- \emph{bằng}, không phải chặn. Do đó
$\|g_S\|_1+|\sigma|=\|g_S\|_1+\|g_{S^c}\|_1=\|g\|_1$, khớp đúng cả vectơ.
Cùng Bước 1:
\[
\|\varphi'(\tau)\|_1\;\le\;\frac{n\|g\|_1}{D(\tau)}\;\le\;\frac{n\beta\|\Gamma\|_1}{L}\;=\;A .
\]
Lấy tích phân trên $[0,1]$ được \eqref{eq:dichchuyen}.

\emph{Bước 5 (làm tròn).} $\quant(u)=u+r(u)$ với $\|r(u)\|_\infty\le\tfrac12$,
nên $\|\quant(u)-\quant(v)\|_1\le\|u-v\|_1+C$; áp vào $u=R(c-\beta\Gamma)$,
$v=R(c)$.
\end{proof}

\begin{hequa}[Chặn dưới cho sai số khôi phục]
\label{hq:chanduoi}
\textsc{[cm]}\;
\[
\|\hat N-N\|_1\;\ge\;
\underbrace{\bigl\|N-\quant(R(c))\bigr\|_1}_{\text{nhãn thật cách tâm bao xa}}
\;-\;(A+C).
\]
Tức là: \textbf{tấn công không thể tốt hơn ``đoán tiên nghiệm lớp của tập phụ
trợ'' quá $A+C$} --- đúng với \emph{mọi} $\epstr$, không cần ngưỡng.
\end{hequa}

\begin{proof}
Bất đẳng thức tam giác đảo áp vào \eqref{eq:bay}.
\end{proof}

\begin{hequa}[Giới hạn]
\label{hq:gioihan}
\textsc{[cm]}\; $\quant\bigl(R(c-\beta\lambda\Gamma)\bigr)\to\quant\bigl(R(c)\bigr)$
khi $\lambda\to0^+$, trừ khi $R(c)$ có tọa độ nửa nguyên. Đây là trường hợp
$A\to0$ của Mệnh đề \ref{md:bay}; huấn luyện đóng đúng vai $\lambda\to0$.
\end{hequa}

\begin{nhanxet}[Thang độ gây ra lỗi; lượng tử hóa khóa lỗi lại]
\label{nx:luongtu}
Dễ quy toàn bộ thất bại cho $\quant$, nhưng như thế là sai địa chỉ:
\begin{itemize}
\item \textbf{$A$ nhỏ, một mình nó, đã làm đầu ra sai.} Theo
  \eqref{eq:dichchuyen}, ngay \emph{trước} khi làm tròn, đầu ra đã bị nhốt
  trong bán kính $A$. Nếu $N$ cách tâm $\Theta(n)$ --- đúng với phân hoạch
  Dirichlet lệch --- thì tấn công đã sai kể cả khi tính bằng số thực chính
  xác tuyệt đối. $\quant$ \textbf{không cần thiết} cho thất bại.
\item \textbf{$\quant$ làm thất bại đó không sửa được.} Bước số thực chỉ
  \emph{nén} tín hiệu, về nguyên tắc còn nhân lại được; $\quant$ cắt mất khả
  năng ấy, vì nhiễu loạn không đẩy nổi một tọa độ qua trọn một đơn vị thì nó
  biến mất vĩnh viễn.
\end{itemize}
Nói gọn: \textbf{thang độ sai gây ra lỗi; lượng tử hóa khóa lỗi lại.} Hệ quả
thực tế: bỏ $\quant$ đi không cứu được ZLG$^+$ hay RLU; chỉ bỏ thước ngoài
mới cứu được.
\end{nhanxet}

\begin{nhanxet}[Cắt không phải chế độ hỏng]
\label{nx:cat}
Tọa độ $x_j<0$ nghĩa là $\beta\Gamma_j>c_j\ge0$: gradient dương lớn trên lớp
$j$, dấu hiệu của lớp \emph{vắng mặt}; cắt về $0$ cho đúng $\hat N_j=0$. Và
theo Bước 1 của chứng minh trên, cắt chỉ làm $D$ tăng, tức nén tín hiệu thêm.
Nên nó \emph{thu nhỏ} quả cầu bẫy chứ không nới --- đi cùng chiều với kết
luận.
\end{nhanxet}

% ===========================================================================
\section{Định lượng: chặn này cắn ở đâu}
\label{sec:dinh-luong}

\begin{nhanxet}[ZLG$^+$: suy giảm liên tục và cắn sớm]
\label{nx:lientuc}
\eqref{eq:bay} nói tỉ lệ khối lượng đầu ra mà \emph{nhãn} được phép điều
khiển nhiều nhất là $(A+C)/n$. Với ZLG$^+$ ($\beta=1$,
$\|c\|_1=nC(1-\epsaux)$), ghép với Bổ đề \ref{bd:dolon}:
\begin{equation}
\frac{A}{n}\;\le\;\frac{2\,w_{\max}\,\epstr}{C\,(1-\epsaux)} .
\label{eq:tilekhoiluong}
\end{equation}
Chú ý $1/C$ ở mẫu số --- đó là lý do chặn cắn \emph{sớm}. Thế các giá trị
minh họa $C=10$, $\epsaux=0{,}25$, $w_{\max}=1$, $n=640$ (đều \textbf{chưa
đo}, đưa vào chỉ để thấy bậc độ lớn):

\begin{center}
\begin{tabular}{@{}ccc@{}}
\toprule
$\epstr$ & $(A+C)/n$ tối đa & khối lượng do tâm quyết định \\
\midrule
$0{,}7$  & $0{,}20$ & $\ge80\%$ \\
$0{,}3$  & $0{,}10$ & $\ge90\%$ \\
$0{,}05$ & $0{,}03$ & $\ge97\%$ \\
\bottomrule
\end{tabular}
\end{center}

Không có ngưỡng và không có sụp đổ đột ngột: suy giảm bắt đầu ngay khi
$\epstr$ còn cao, rồi sâu dần. Đây là chỗ báo cáo khớp với quan sát rằng
ZLG$^+$ đã giảm từ lúc mô hình còn kém.

\eqref{eq:tilekhoiluong} cho một \textbf{tiên đoán phân biệt được}: vì
$A/n\propto1/C$, ở cùng $\epstr$ thì CIFAR100 phải suy giảm nặng hơn CIFAR10
khoảng một bậc.
\end{nhanxet}

\begin{nhanxet}[RLU: chặn yếu, và ở sai số cao thì rỗng]
\label{nx:rluyeu}
\textsc{[cm]}\; Với RLU, $\beta=1/n$ và $\|c\|_1=\bar\rho$, nên
$A/n\le2\epstr/\bar\rho$ --- \emph{không} có thừa số $1/C$. Với
$\bar\rho\approx1{,}0988$ (Mệnh đề \ref{md:ghim}.2), độ dốc $n/\|c\|_1$ của
hai phương pháp lệch nhau bốn bậc:

\begin{center}
\begin{tabular}{@{}lccc@{}}
\toprule
 & độ dốc $n/\|c\|_1$ & $A/n$ & tại $\epstr=0{,}7$ \\
\midrule
ZLG$^+$ & $1/\bigl(C(1-\epsaux)\bigr)\approx0{,}133$ & $\le0{,}267\,\epstr$ & $0{,}19$ \\
RLU     & $n/\bar\rho\approx583$ & $\le1{,}82\,\epstr$ & $1{,}27$ \\
\bottomrule
\end{tabular}
\end{center}

Chặn này \emph{không} rỗng, nhưng yếu. Từ một tâm gần đều, điểm xa nhất trên
đơn hình cách $2n(1-1/C)=1{,}8n$, nên bán kính $A$ vẫn loại trừ được một phần
dải khả dĩ:

\begin{center}
\begin{tabular}{@{}ccc@{}}
\toprule
$\epstr$ & ZLG$^+$ loại trừ & RLU loại trừ \\
\midrule
$0{,}7$  & $90\%$ & $29\%$ \\
$0{,}5$  & $93\%$ & $49\%$ \\
$0{,}3$  & $96\%$ & $70\%$ \\
$0{,}1$  & $99\%$ & $90\%$ \\
\bottomrule
\end{tabular}
\end{center}

Phát biểu đúng: chặn của RLU \textbf{yếu hơn ZLG$^+$ đúng $6{,}8$ lần}, bằng
tỉ số hai độ dốc $C(1-\epsaux)/\bar\rho$. RLU phải xuống tới
$\epstr\approx0{,}1$ mới đạt mức đảm bảo mà ZLG$^+$ đã có từ $0{,}7$.

Gốc rễ của chênh lệch: ZLG$^+$ so tín hiệu với một thước gộp $C$ lớp, còn RLU
so với một phân phối xác suất đã chuẩn hóa ($\|c\|_1=\bar\rho\approx1{,}1$).
Thiếu đúng thừa số $C$.

Một đường thoát khả dĩ, \textbf{chưa kiểm}: $\|\Gamma\|_1\le2n\epstr$ là chặn
xấu nhất, bỏ qua triệt tiêu giữa $a$ và $Q^{\!\top}a$ trong Bổ đề
\ref{bd:tinhieu}. Nếu $\|\Gamma\|_1$ thực tế nhỏ hơn chặn đó vài lần thì RLU
cũng rơi vào vùng chặn cắn.
\end{nhanxet}

% ===========================================================================
\section{Kết luận}
\label{sec:ket-luan}

\subsection{Chuỗi lập luận, gói lại}

\begin{enumerate}
\item Cả ZLG$^+$ và RLU đều quy về $\hat N=\quant\bigl(R(c-\beta\Gamma)\bigr)$
  với $c$ lấy hoàn toàn từ tập phụ trợ (Mệnh đề \ref{md:dangchuan}).
\item $\|\Gamma\|_1\le2n\epstr$: độ lớn gradient lớp cuối \emph{chính là} sai
  số huấn luyện, không có đường thoát (Bổ đề \ref{bd:dolon}).
\item $\|c\|_1$ bị ghim: của ZLG$^+$ \emph{tăng} về $nC$, của RLU là hằng số
  không đụng tới dữ liệu (Mệnh đề \ref{md:ghim}).
\item Nên $A=\Theta(\epstr)$ (Hệ quả \ref{hq:Aeps}), và $A$ chặn trực tiếp
  khoảng cách từ đầu ra tới tâm: $\|\hat N-\text{tâm}\|_1\le A+C$ (Mệnh đề
  \ref{md:bay}).
\item Tâm không chứa thông tin nhãn (Nhận xét \ref{nx:tam}).
\end{enumerate}

\begin{quote}
\textbf{Khi mô hình được huấn luyện tốt, ZLG$^+$ và RLU không còn đoán nhãn
của client --- chúng đoán phân bố lớp của chính dữ liệu phụ trợ của mình.}
\end{quote}

Hay ngắn hơn: \emph{suy giảm không phải vì tín hiệu yếu đi, mà vì nó được đo
bằng một cái thước không chịu yếu theo.}

\subsection{Năm chỗ kết luận này không phủ}

\begin{enumerate}
\item\label{item:hvong} \textbf{Với RLU, điều đã chứng minh không phải là
  đầu ra.} Theo Nhận xét \ref{nx:haivong}, RLU chạy dạng chuẩn hai vòng và
  trả về kết quả vòng hai, mà tâm vòng hai lại được dựng từ ước lượng vòng
  một --- nên nó \emph{phụ thuộc nhãn}. Đây là chỗ yếu nghiêm trọng nhất của
  báo cáo: mắt xích ``tâm vô can'' hỏng đúng ở bước sinh ra đầu ra. Đóng lại
  cần một lập luận hợp thành (tâm vòng hai là hàm của $N_1$, mà $N_1$ đã bị
  nhốt), nhưng vòng tinh chỉnh chạy tới $5$ lần, có điều kiện dừng, và độ
  nhạy của nó theo $N$ \textbf{chưa đo}.

\item \textbf{Với RLU, chặn còn yếu hơn ZLG$^+$ $6{,}8$ lần} (Nhận xét
  \ref{nx:rluyeu}), nên ngay ở vòng một kết luận cũng chỉ chặt cho ZLG$^+$.
  Chỉ đóng được bằng số đo $\|\Gamma\|_1$ thật, không đóng được bằng thêm đại
  số.
\item \textbf{Chặn chỉ cắn khi client lệch khỏi tâm.} Nếu client tình cờ có
  phân bố gần tâm thì Hệ quả \ref{hq:chanduoi} không nói gì, và hai phương
  pháp sẽ \emph{trông như} vẫn chạy tốt dù không suy ra gì từ nhãn. Tiên đoán
  kèm theo: hiệu năng còn lại phải tương quan với độ lệch của client khỏi
  phân bố đều, chứ không với chất lượng suy luận.
\item \textbf{Tâm phụ thuộc phụ trợ.} \emph{Việc} bị nhốt quanh một vectơ
  không phụ thuộc nhãn là chứng minh được; \emph{vectơ nào} thì do tập phụ
  trợ quyết định.
\item \textbf{Huấn luyện không chỉ làm co $\Gamma$, nó còn đổi hướng.} $Q$
  trôi theo huấn luyện. Báo cáo này chỉ xử lý thành phần độ lớn --- thành
  phần trội, nhưng không phải thành phần duy nhất.
\end{enumerate}

\subsection{Những gì cần đo}

\begin{enumerate}
\item \textbf{Kiểm ràng buộc cứng \eqref{eq:bay}.} Đo $A$ và
  $\|\hat N-\quant(R(c))\|_1$ theo từng cửa sổ; bất đẳng thức phải đúng ở
  \emph{mọi} cửa sổ. Một vi phạm duy nhất bác bỏ dạng chuẩn
  \eqref{eq:dangchuan}. Đây là phép thử mạnh nhất vì nó nhị phân, không phải
  tương quan.
\item \textbf{Độ nhạy của \code{estimate\_static\_RLU\_with\_posterior}
  theo $N$.} Nếu tâm vòng hai dịch ít hơn hẳn so với $N_1$ thì lập luận hợp
  thành đóng được và hồ sơ RLU hoàn chỉnh (chỗ không phủ số
  \ref{item:hvong}). Ưu tiên cao nhất cho RLU.

\item \textbf{$\|\Gamma\|_1$ thật so với chặn $2n\epstr$.} Quyết định liệu
  Mệnh đề \ref{md:bay} có cắn được cho RLU ở $\epstr$ cao hay không. Nếu
  $\|\Gamma\|_1/(2n\epstr)<\bar\rho/2$ thì hồ sơ RLU đóng lại.
\item \textbf{Tiên đoán $1/C$}: chạy cùng $\epstr$ trên CIFAR10 và CIFAR100,
  kiểm $A/n$ của ZLG$^+$ có giảm khoảng một bậc không. RLU \emph{không} được
  giảm theo cách đó; nếu cả hai giảm như nhau thì $1/C$ không phải cơ chế
  thật.
\item \textbf{$w_{\max}=\max_i s_i/\bar O$} --- giả thiết duy nhất còn lại
  của Bổ đề \ref{bd:dolon} cho $\tilde G$.
\item \textbf{Hành vi của $\bar O$ khi huấn luyện} (Nhận xét \ref{nx:Obar}).
\item \textbf{Độ lệch của client}: tương quan giữa irec còn lại và
  $\|N-\quant(R(c))\|_1$ (chỗ không phủ số 2).
\end{enumerate}

% ===========================================================================
\appendix
\section{Vị trí trong mã nguồn của mọi khẳng định \textsc{[cm]}}
\label{sec:phuluc}

\begin{center}
\begin{tabular}{@{}lll@{}}
\toprule
\textbf{Khẳng định} & \textbf{Tệp} & \textbf{Hàm} \\
\midrule
ZLG$^+$ có dạng \eqref{eq:dangchuan} & \code{client/attacks.py} & \code{ZLGp\_attack} \\
RLU có dạng \eqref{eq:dangchuan} & \code{client/client\_utils.py} & \code{estimated\_entropy\_from\_grad} \\
$c_{\text{RLU}}=\bar\rho\,\mathrm{softmax}(\bar\mu)$ & \code{client/attacks.py} & \code{RLU\_attack} \\
$\bar\rho$ không phụ thuộc mô hình & \code{client/attacks.py} & \code{rlu\_rho\_mean} \\
$\bar p_j=1-\bar\varepsilon_{\mathrm{aux},j}$ & \code{client/client\_utils.py} & \code{estimate\_static\_ZLG} \\
$\bar O=\mathrm{mean}(s_i)$ trên phụ trợ & \code{client/client\_utils.py} & \code{estimate\_static\_ZLG} \\
Phụ trợ lấy từ tập test & \code{main.py} & dòng \code{aux\_dataset = ...} \\
\bottomrule
\end{tabular}
\end{center}

\end{document}
